\documentclass[11pt,a4paper]{article}

\usepackage[T1]{fontenc}
\usepackage[utf8]{inputenc}
\usepackage[polish]{babel}
\usepackage{geometry}
\usepackage{hyperref}
\usepackage{enumitem}
\usepackage{longtable}
\usepackage{array}
\usepackage{xcolor}
\usepackage{listings}
\usepackage{titlesec}

\geometry{margin=2.5cm}

\hypersetup{
    colorlinks=true,
    linkcolor=blue,
    urlcolor=blue
}

\lstset{
    basicstyle=\ttfamily\small,
    breaklines=true,
    frame=single,
    columns=fullflexible
}

\titleformat{\section}{\large\bfseries}{\thesection.}{0.5em}{}
\titleformat{\subsection}{\normalsize\bfseries}{\thesubsection.}{0.5em}{}

\title{Backend aplikacji -- dokumentacja developerska}
\author{}
\date{}

\begin{document}

\maketitle

\section{Cel dokumentu}

Ten dokument opisuje aktualny backend aplikacji, sposób połączenia modułów, główne przepływy danych oraz sposób uruchomienia i testowania projektu lokalnie. Ma służyć kolejnym programistom jako szybkie wejście do kodu, tak aby mogli bezpiecznie rozwijać backend i podłączać frontend.

\section{Stack technologiczny}

\begin{itemize}[leftmargin=1.2cm]
    \item Python
    \item Flask
    \item Flask-SQLAlchemy
    \item Flask-JWT-Extended
    \item Marshmallow
    \item SQLite na potrzeby lokalnego developmentu
    \item pytest do testów
\end{itemize}

Aktualnie projekt działa na \textbf{SQLite}, co upraszcza lokalne uruchomienie. Jest to rozwiązanie developerskie. Docelowo architektura nadal nadaje się do późniejszego powrotu do PostgreSQL.

\section{Struktura katalogów}

Najważniejsze elementy backendu:

\begin{lstlisting}
backend/
|-- app/
|   |-- __init__.py
|   |-- config.py
|   |-- extensions.py
|   |-- seed_data.py
|   |-- api/
|   |-- models/
|   |-- repositories/
|   |-- schemas/
|   |-- services/
|   |-- utils/
|   `-- cli/
|-- tests/
|-- requirements.txt
`-- run.py
\end{lstlisting}

\subsection{Odpowiedzialność warstw}

\subsubsection*{app/api/}
Warstwa HTTP. Zawiera blueprinty Flask i endpointy REST. Odpowiada za:
\begin{itemize}[leftmargin=1.2cm]
    \item odbiór requestu,
    \item walidację danych wejściowych przez schemat,
    \item wywołanie odpowiedniego serwisu,
    \item zwrócenie spójnej odpowiedzi JSON.
\end{itemize}

Route'y powinny być cienkie. Logika biznesowa nie powinna trafiać tutaj.

\subsubsection*{app/services/}
Główna logika biznesowa systemu. Tu realizowane są:
\begin{itemize}[leftmargin=1.2cm]
    \item rejestracja i logowanie,
    \item wyszukiwanie usług i terapeutów,
    \item liczenie dostępności,
    \item rezerwacja i anulowanie wizyt,
    \item obsługa opinii,
    \item operacje admin/staff.
\end{itemize}

To najważniejsza warstwa do dalszej rozbudowy.

\subsubsection*{app/repositories/}
Warstwa dostępu do danych. Zawiera zapytania do modeli i bazy. W projekcie repozytoria są używane częściowo --- część serwisów korzysta też bezpośrednio z modeli ORM.

\subsubsection*{app/models/}
Modele SQLAlchemy, czyli definicje tabel i relacji.

\subsubsection*{app/schemas/}
Schematy Marshmallow do:
\begin{itemize}[leftmargin=1.2cm]
    \item walidacji requestów,
    \item serializacji odpowiedzi.
\end{itemize}

\subsubsection*{app/utils/}
Wspólne helpery:
\begin{itemize}[leftmargin=1.2cm]
    \item błędy,
    \item role,
    \item paginacja,
    \item funkcje czasu,
    \item auth helpers.
\end{itemize}

\subsubsection*{app/seed\_data.py}
Dane startowe do developmentu i testów.

\subsubsection*{tests/}
Testy backendu.

\section{Główne moduły backendu}

\subsection{Auth}
Pliki:
\begin{itemize}[leftmargin=1.2cm]
    \item \texttt{app/api/auth\_routes.py}
    \item \texttt{app/services/auth\_service.py}
\end{itemize}

Obsługuje:
\begin{itemize}[leftmargin=1.2cm]
    \item rejestrację,
    \item logowanie,
    \item refresh token,
    \item logout,
    \item pobieranie aktualnego użytkownika.
\end{itemize}

\subsection{Katalog usług i terapeutów}
Pliki:
\begin{itemize}[leftmargin=1.2cm]
    \item \texttt{app/api/services\_routes.py}
    \item \texttt{app/api/therapists\_routes.py}
    \item \texttt{app/services/catalog\_service.py}
\end{itemize}

Obsługuje:
\begin{itemize}[leftmargin=1.2cm]
    \item listę usług,
    \item szczegóły usługi,
    \item terapeutów przypisanych do usługi,
    \item publiczny profil terapeuty.
\end{itemize}

\subsection{Dostępność}
Pliki:
\begin{itemize}[leftmargin=1.2cm]
    \item \texttt{app/api/availability\_routes.py}
    \item \texttt{app/services/availability\_service.py}
\end{itemize}

Obsługuje liczenie wolnych terminów na podstawie:
\begin{itemize}[leftmargin=1.2cm]
    \item reguł grafiku,
    \item wyjątków,
    \item istniejących wizyt,
    \item czasu trwania usługi.
\end{itemize}

Sloty nie są trzymane w tabeli. Są generowane dynamicznie.

\subsection{Wizyty}
Pliki:
\begin{itemize}[leftmargin=1.2cm]
    \item \texttt{app/api/appointments\_routes.py}
    \item \texttt{app/services/appointment\_service.py}
\end{itemize}

Obsługuje:
\begin{itemize}[leftmargin=1.2cm]
    \item utworzenie rezerwacji,
    \item listę wizyt pacjenta,
    \item szczegóły wizyty,
    \item anulowanie wizyty.
\end{itemize}

\subsection{Historia konsultacji}
Pliki:
\begin{itemize}[leftmargin=1.2cm]
    \item \texttt{app/api/consultations\_routes.py}
    \item \texttt{app/services/consultation\_service.py}
\end{itemize}

Obsługuje:
\begin{itemize}[leftmargin=1.2cm]
    \item historię zakończonych wizyt pacjenta,
    \item opis konsultacji.
\end{itemize}

\subsection{Opinie}
Pliki:
\begin{itemize}[leftmargin=1.2cm]
    \item \texttt{app/api/reviews\_routes.py}
    \item \texttt{app/services/review\_service.py}
\end{itemize}

Obsługuje:
\begin{itemize}[leftmargin=1.2cm]
    \item publiczne opinie terapeuty,
    \item dodanie opinii po zakończonej wizycie,
    \item przeliczanie średniej oceny terapeuty.
\end{itemize}

\subsection{Admin}
Pliki:
\begin{itemize}[leftmargin=1.2cm]
    \item \texttt{app/api/admin\_services\_routes.py}
    \item \texttt{app/api/admin\_therapists\_routes.py}
    \item \texttt{app/api/admin\_availability\_routes.py}
    \item \texttt{app/services/admin\_service.py}
    \item \texttt{app/services/therapist\_admin\_service.py}
    \item \texttt{app/services/schedule\_admin\_service.py}
\end{itemize}

Obsługuje:
\begin{itemize}[leftmargin=1.2cm]
    \item zarządzanie usługami,
    \item tworzenie i edycję terapeutów,
    \item przypisywanie usług do terapeutów,
    \item zarządzanie grafikami i wyjątkami.
\end{itemize}

\subsection{Staff}
Pliki:
\begin{itemize}[leftmargin=1.2cm]
    \item \texttt{app/api/staff\_appointments\_routes.py}
    \item \texttt{app/services/staff\_appointment\_service.py}
\end{itemize}

Obsługuje:
\begin{itemize}[leftmargin=1.2cm]
    \item listę wizyt do obsługi,
    \item zmianę statusów wizyt,
    \item dodanie lub edycję opisu konsultacji.
\end{itemize}

\section{Modele i relacje}

Najważniejsze encje:
\begin{itemize}[leftmargin=1.2cm]
    \item \texttt{User}
    \item \texttt{Role}
    \item \texttt{UserRole}
    \item \texttt{TherapistProfile}
    \item \texttt{Service}
    \item \texttt{TherapistService}
    \item \texttt{AvailabilityRule}
    \item \texttt{AvailabilityException}
    \item \texttt{Appointment}
    \item \texttt{ConsultationSummary}
    \item \texttt{Review}
    \item \texttt{EmailNotification}
\end{itemize}

\subsection{Najważniejsze relacje}

\begin{itemize}[leftmargin=1.2cm]
    \item użytkownik może mieć wiele ról,
    \item terapeuta jest powiązany z użytkownikiem 1:1,
    \item terapeuta może realizować wiele usług,
    \item usługa może być realizowana przez wielu terapeutów,
    \item terapeuta ma wiele reguł grafiku i wyjątków,
    \item pacjent ma wiele wizyt,
    \item wizyta należy do jednego pacjenta, jednej usługi i jednego terapeuty,
    \item wizyta może mieć jedno podsumowanie konsultacji,
    \item wizyta może mieć jedną opinię.
\end{itemize}

\subsection{Snapshoty w wizytach}

W tabeli \texttt{Appointment} przechowywane są snapshoty:
\begin{itemize}[leftmargin=1.2cm]
    \item nazwy usługi,
    \item opisu usługi,
    \item długości wizyty,
    \item ceny,
    \item imienia i nazwiska terapeuty,
    \item tytułu terapeuty.
\end{itemize}

Dzięki temu historia wizyt nie zmienia się po edycji katalogu usług lub profilu terapeuty.

\section{Role i uprawnienia}

Role:
\begin{itemize}[leftmargin=1.2cm]
    \item \texttt{PATIENT}
    \item \texttt{THERAPIST}
    \item \texttt{ADMIN}
\end{itemize}

\subsection{Pacjent}
Może:
\begin{itemize}[leftmargin=1.2cm]
    \item się zarejestrować,
    \item zalogować,
    \item przeglądać usługi i terapeutów,
    \item sprawdzać dostępność,
    \item rezerwować wizytę,
    \item anulować własną wizytę,
    \item przeglądać własne wizyty i konsultacje,
    \item dodawać opinię po zakończonej wizycie.
\end{itemize}

\subsection{Terapeuta}
Może:
\begin{itemize}[leftmargin=1.2cm]
    \item przeglądać staffowe wizyty przypisane do siebie,
    \item zmieniać statusy swoich wizyt,
    \item dodawać opis konsultacji do obsługiwanych wizyt.
\end{itemize}

\subsection{Admin}
Może:
\begin{itemize}[leftmargin=1.2cm]
    \item zarządzać usługami,
    \item zarządzać terapeutami,
    \item przypisywać usługi,
    \item zarządzać grafikami,
    \item przeglądać i edytować staffowe wizyty.
\end{itemize}

\section{Przepływ requestu}

Typowy przepływ wygląda tak:
\begin{enumerate}[leftmargin=1.2cm]
    \item request trafia do endpointu Flask w \texttt{app/api/...},
    \item dane wejściowe są walidowane przez schemat w \texttt{app/schemas/...},
    \item route wywołuje serwis z \texttt{app/services/...},
    \item serwis wykonuje logikę biznesową,
    \item w razie potrzeby serwis korzysta z repozytorium lub modeli SQLAlchemy,
    \item wynik wraca do route,
    \item odpowiedź jest opakowana do formatu:
    \begin{itemize}
        \item \texttt{data},
        \item opcjonalnie \texttt{meta},
        \item albo \texttt{error}.
    \end{itemize}
\end{enumerate}

\subsection{Przykład: POST /api/v1/appointments}

\begin{itemize}[leftmargin=1.2cm]
    \item route: \texttt{appointments\_routes.py}
    \item schema: \texttt{AppointmentCreateSchema}
    \item service: \texttt{AppointmentService.create(...)}
\end{itemize}

Serwis:
\begin{itemize}[leftmargin=1.2cm]
    \item sprawdza usługę,
    \item sprawdza terapeutę,
    \item sprawdza powiązanie terapeuta-usługa,
    \item sprawdza, czy termin istnieje w dostępności,
    \item sprawdza konflikty,
    \item zapisuje wizytę,
    \item tworzy powiadomienia email.
\end{itemize}

\section{Format odpowiedzi API}

\subsection{Sukces}
\begin{lstlisting}
{
  "data": {}
}
\end{lstlisting}

\subsection{Sukces z paginacją}
\begin{lstlisting}
{
  "data": [],
  "meta": {
    "page": 1,
    "pageSize": 10,
    "total": 100
  }
}
\end{lstlisting}

\subsection{Błąd}
\begin{lstlisting}
{
  "error": {
    "code": "APPOINTMENT_SLOT_NOT_AVAILABLE",
    "message": "Wybrany termin nie jest dostępny."
  }
}
\end{lstlisting}

Walidacja może zwracać też \texttt{details}.

\section{Dostępne endpointy}

\subsection{Public / patient}

\subsubsection*{Auth}
\begin{itemize}[leftmargin=1.2cm]
    \item \texttt{POST /api/v1/auth/register}
    \item \texttt{POST /api/v1/auth/login}
    \item \texttt{POST /api/v1/auth/refresh}
    \item \texttt{POST /api/v1/auth/logout}
    \item \texttt{GET /api/v1/users/me}
\end{itemize}

\subsubsection*{Katalog}
\begin{itemize}[leftmargin=1.2cm]
    \item \texttt{GET /api/v1/services}
    \item \texttt{GET /api/v1/services/<service\_id>}
    \item \texttt{GET /api/v1/services/<service\_id>/therapists}
    \item \texttt{GET /api/v1/therapists/<therapist\_id>}
    \item \texttt{GET /api/v1/therapists/<therapist\_id>/reviews}
\end{itemize}

\subsubsection*{Dostępność}
\begin{itemize}[leftmargin=1.2cm]
    \item \texttt{GET /api/v1/availability}
\end{itemize}

\subsubsection*{Wizyty}
\begin{itemize}[leftmargin=1.2cm]
    \item \texttt{POST /api/v1/appointments}
    \item \texttt{GET /api/v1/appointments/me}
    \item \texttt{GET /api/v1/appointments/<appointment\_id>}
    \item \texttt{POST /api/v1/appointments/<appointment\_id>/cancel}
\end{itemize}

\subsubsection*{Konsultacje i opinie}
\begin{itemize}[leftmargin=1.2cm]
    \item \texttt{GET /api/v1/consultations/me}
    \item \texttt{POST /api/v1/appointments/<appointment\_id>/review}
\end{itemize}

\subsection{Admin}
\begin{itemize}[leftmargin=1.2cm]
    \item \texttt{GET /api/v1/admin/services}
    \item \texttt{POST /api/v1/admin/services}
    \item \texttt{PATCH /api/v1/admin/services/<service\_id>}
    \item \texttt{GET /api/v1/admin/therapists}
    \item \texttt{POST /api/v1/admin/therapists}
    \item \texttt{PATCH /api/v1/admin/therapists/<therapist\_id>}
    \item \texttt{POST /api/v1/admin/therapists/<therapist\_id>/services}
    \item \texttt{DELETE /api/v1/admin/therapists/<therapist\_id>/services/<service\_id>}
    \item \texttt{GET /api/v1/admin/therapists/<therapist\_id>/availability-rules}
    \item \texttt{POST /api/v1/admin/therapists/<therapist\_id>/availability-rules}
    \item \texttt{PATCH /api/v1/admin/availability-rules/<rule\_id>}
    \item \texttt{DELETE /api/v1/admin/availability-rules/<rule\_id>}
    \item \texttt{GET /api/v1/admin/therapists/<therapist\_id>/availability-exceptions}
    \item \texttt{POST /api/v1/admin/therapists/<therapist\_id>/availability-exceptions}
\end{itemize}

\subsection{Staff}
\begin{itemize}[leftmargin=1.2cm]
    \item \texttt{GET /api/v1/staff/appointments}
    \item \texttt{PATCH /api/v1/staff/appointments/<appointment\_id>/status}
    \item \texttt{PUT /api/v1/staff/appointments/<appointment\_id>/consultation-summary}
\end{itemize}

\subsection{Techniczne}
\begin{itemize}[leftmargin=1.2cm]
    \item \texttt{GET /api/v1/health}
\end{itemize}

\section{Jak uruchomić projekt lokalnie}

\subsection{1. Wejdź do katalogu backend}
\begin{lstlisting}
cd backend
\end{lstlisting}

\subsection{2. Utwórz i aktywuj virtualenv}
macOS / Linux:
\begin{lstlisting}
python3 -m venv .venv
source .venv/bin/activate
\end{lstlisting}

\subsection{3. Zainstaluj zależności}
\begin{lstlisting}
pip install -r requirements.txt
\end{lstlisting}

\subsection{4. Ustaw zmienne środowiskowe}

Utwórz plik \texttt{.env} albo ustaw zmienne w shellu.

Przykład:
\begin{lstlisting}
FLASK_ENV=development
DATABASE_URL=sqlite:///app.db
SECRET_KEY=change-me
JWT_SECRET_KEY=change-me-too
CORS_ORIGINS=http://localhost:3000
APP_TIMEZONE=Europe/Warsaw
MAIL_SENDER=noreply@example.com
MAIL_BACKEND=console
\end{lstlisting}

\subsection{5. Usuń starą bazę, jeśli chcesz startować od zera}
\begin{lstlisting}
rm -f app.db
rm -f instance/app.db
\end{lstlisting}

\subsection{6. Uruchom backend}
\begin{lstlisting}
python run.py
\end{lstlisting}

\subsubsection*{Co robi run.py}
\begin{itemize}[leftmargin=1.2cm]
    \item tworzy aplikację Flask,
    \item wykonuje \texttt{db.create\_all()},
    \item uruchamia \texttt{seed\_database()},
    \item startuje serwer developerski.
\end{itemize}

To rozwiązanie jest wygodne lokalnie. Na późniejszym etapie warto przejść na migracje.

\section{Dane testowe seed}

Przykładowe konta:
\begin{itemize}[leftmargin=1.2cm]
    \item admin: \texttt{admin@example.com} / \texttt{Admin123!}
    \item terapeuta: \texttt{anna@example.com} / \texttt{Password123!}
    \item terapeuta: \texttt{piotr@example.com} / \texttt{Password123!}
    \item pacjent: \texttt{jan@example.com} / \texttt{Password123!}
    \item pacjent: \texttt{ola@example.com} / \texttt{Password123!}
\end{itemize}

W seedzie są też:
\begin{itemize}[leftmargin=1.2cm]
    \item usługi,
    \item profile terapeutów,
    \item przypisania usług,
    \item grafiki,
    \item przykładowe wizyty,
    \item przykładowa opinia,
    \item powiadomienia email.
\end{itemize}

\section{Jak sprawdzić, czy backend działa}

\subsection{1. Healthcheck}
Wejdź w przeglądarce albo użyj curl:

\begin{lstlisting}
curl http://127.0.0.1:5000/api/v1/health
\end{lstlisting}

Oczekiwany wynik:
\begin{lstlisting}
{
  "data": {
    "status": "ok"
  }
}
\end{lstlisting}

\subsection{2. Lista usług}
\begin{lstlisting}
curl http://127.0.0.1:5000/api/v1/services
\end{lstlisting}

Powinieneś dostać listę usług.

\subsection{3. Logowanie}
\begin{lstlisting}
curl -X POST http://127.0.0.1:5000/api/v1/auth/login \
  -H "Content-Type: application/json" \
  -d '{
    "email": "jan@example.com",
    "password": "Password123!"
  }'
\end{lstlisting}

W odpowiedzi dostaniesz token.

\subsection{4. Autoryzowany request}
Podstaw \texttt{accessToken} do:
\begin{lstlisting}
curl http://127.0.0.1:5000/api/v1/users/me \
  -H "Authorization: Bearer <accessToken>"
\end{lstlisting}

\section{Przykładowy manualny test przepływu pacjenta}

\begin{enumerate}[leftmargin=1.2cm]
    \item \texttt{POST /auth/login}
    \item \texttt{GET /services}
    \item pobierz \texttt{service\_id}
    \item \texttt{GET /services/<service\_id>/therapists}
    \item pobierz \texttt{therapist\_id}
    \item \texttt{GET /availability?serviceId=...\&therapistId=...\&from=...\&to=...}
    \item wybierz \texttt{startAt}
    \item \texttt{POST /appointments}
    \item \texttt{GET /appointments/me}
    \item \texttt{POST /appointments/<id>/cancel}
\end{enumerate}

To jest podstawowy flow, który powinien działać przed podpinaniem frontendu.

\section{Jak uruchomić testy}

W katalogu backend:

\begin{lstlisting}
pytest tests -q
\end{lstlisting}

\subsection{Co testujemy}
Testy obejmują:
\begin{itemize}[leftmargin=1.2cm]
    \item auth,
    \item katalog,
    \item dostępność,
    \item wizyty,
    \item review/admin/staff.
\end{itemize}

\subsection{Jak działają testy}
\begin{itemize}[leftmargin=1.2cm]
    \item testy używają konfiguracji \texttt{testing},
    \item baza jest tworzona w pamięci,
    \item dane seedujące są ładowane automatycznie w fixture.
\end{itemize}

\section{Najważniejsze miejsca do dalszej rozbudowy}

\subsection{Backend}
Najczęstsze miejsca rozwoju:
\begin{itemize}[leftmargin=1.2cm]
    \item nowe endpointy w \texttt{app/api/},
    \item nowa logika w \texttt{app/services/},
    \item nowe walidacje w \texttt{app/schemas/},
    \item rozszerzenia modeli w \texttt{app/models/},
    \item dodatkowe testy w \texttt{tests/}.
\end{itemize}

\subsection{Frontend}
Frontend powinien korzystać głównie z tych endpointów:
\begin{itemize}[leftmargin=1.2cm]
    \item auth,
    \item services,
    \item therapists,
    \item availability,
    \item appointments,
    \item consultations,
    \item reviews.
\end{itemize}

Najpierw warto podłączyć:
\begin{enumerate}[leftmargin=1.2cm]
    \item login/register,
    \item listę usług,
    \item terapeutów,
    \item dostępność,
    \item rezerwację,
    \item moje wizyty.
\end{enumerate}

\section{Ważne uwagi techniczne}

\subsection{1. SQLite jest tymczasowe}
Dobre do developmentu lokalnego, ale nie jest docelowym rozwiązaniem dla rezerwacji terminów na większą skalę.

\subsection{2. db.create\_all() jest developerskie}
Obecnie wygodne, ale później powinno zostać zastąpione przez migracje.

\subsection{3. Logika biznesowa powinna trafiać do serwisów}
Nie rozbudowuj route'ów o skomplikowaną logikę.

\subsection{4. Seed jest częścią rozwoju lokalnego}
Jeśli dodajesz nowe moduły, zadbaj o aktualizację \texttt{seed\_data.py}, żeby inni mogli łatwo uruchomić projekt.

\subsection{5. Testy trzeba rozwijać razem z backendem}
Każdy nowy endpoint biznesowy powinien dostać test.

\section{Rekomendowany sposób pracy nad backendem}

Przy dodawaniu nowej funkcjonalności najlepiej iść w tej kolejności:
\begin{enumerate}[leftmargin=1.2cm]
    \item doprecyzować model danych,
    \item dodać lub poprawić schemat wejścia/wyjścia,
    \item zaimplementować logikę w serwisie,
    \item wystawić endpoint w \texttt{api/},
    \item dodać testy,
    \item zaktualizować seed, jeśli potrzeba.
\end{enumerate}

To pozwala zachować porządek i przewidywalność kodu.

\section{Krótka lista startowa dla nowego programisty}

Jeśli wchodzisz do projektu po raz pierwszy:
\begin{enumerate}[leftmargin=1.2cm]
    \item sklonuj repo,
    \item przejdź do \texttt{backend/},
    \item utwórz virtualenv,
    \item zainstaluj zależności,
    \item ustaw \texttt{.env},
    \item usuń starą bazę SQLite,
    \item uruchom \texttt{python run.py},
    \item sprawdź \texttt{/api/v1/health},
    \item zaloguj się kontem z seeda,
    \item przejdź podstawowy flow pacjenta,
    \item uruchom \texttt{pytest tests -q}.
\end{enumerate}

Po tych krokach powinieneś mieć działające lokalne środowisko i rozumieć podstawowy przepływ aplikacji.

\end{document}